% This file contains the detailed methods for the nf-core/viralmetagenome manuscript

\section*{Data Availability}
The nf-core/viralmetagenome pipeline is freely available at \url{https://github.com/nf-core/viralmetagenome}.
Scripts used to generate the results as well as parts of the results are available on github at \url{https://github.com/Joon-Klaps/nf-core-viralmetagenome-manuscript}.

\section*{Supplementary Methods}

\subsection*{S1. Evaluation of nf-core/viralmetagenome on public data}
A metagenomic dataset was constructed from publically available data by selecting a random subset of 2,000 records from the NCBI Virus database, targeting the following viruses: \textit{Orthonairovirus haemorrhagiae}, \textit{Mammarenavirus lassaense}, \textit{Zika virus}, \textit{West Nile virus}, \textit{Monkeypox virus}, \textit{Influenza A virus}, \textit{Severe acute respiratory syndrome coronavirus 2}, and \textit{Human respiratory syncytial virus A}. Records were filtered to retain only those with associated SRA data, sequenced on the Illumina platform, and with a LibraryStrategy other than `Amplicon'. From the resulting pool, 28 samples were randomly selected and downloaded using nf-core/fetchngs. Samples were then run with nf-core/viralmetagenome using default options with the exception of \texttt{--cluster\_method 'mmseqs-cluster'} due to sequence size limitations of cdhit, and \texttt{--skip\_read\_classification} to speed up the analysis.

For plant pathogens, two samples were randomly selected from the SRA at every query: \textit{Tobacco mosaic virus}, \textit{Tomato spotted wilt virus}, \textit{Cucumber mosaic virus}, and \textit{Potato virus Y}. These were also downloaded using nf-core/fetchngs and processed with viralmetagenome using the same parameters (\texttt{--cluster\_method 'mmseqs-cluster'}, \texttt{--skip\_read\_classification}), with the addition of the Virosaurus plant virus database as annotation database (\texttt{--annotation\_db} \url{https://ftp.expasy.org/databases/viralzone/2020_4/virosaurus98_plant-20200330.fas.gz}) and RVDB's v29.0 unclustered database as reference pool (\texttt{--reference\_pool} \url{https://rvdb.dbi.udel.edu/download/U-RVDBv29.0.fasta.gz}).

\subsection*{S2. Sample mixture resolution}
To evaluate the pipeline's ability to distinguish viral genomes from the same species in mixed samples, we simulated a set of HIV-1 co-infections. We retrieved a random subset of 2,000 HIV-1 genomes from NCBI Virus (taxid: 11676) and selected MN090277.1 as the constant representative genome. We then calculated the pairwise Average Nucleotide Identity (ANI) between all 2,000 genomes and the representative using MMseqs2 \citep{Steinegger2017-ci}. Based on these ANI values, we binned the genomes into 0.5\% intervals (ranging from 74\% to 100\%) and selected one target genome from each bin (n=38). Synthetic paired-end reads were generated for both the selected target genome and the representative genome using InSilicoSeq v2.0.1 \citep{Gourle2019-ox} in 'kde' mode with the 'MiSeq' error model. Mixtures were simulated at a 50:50 abundance ratio with 50x coverage for each genome.

The simulated mixtures were analyzed using viralmetagenome with the parameters \texttt{--skip\_preprocessing}, \texttt{--skip\_read\_classification}, and with U-RVDBv29.0 as a reference pool. We characterised the pipeline's performance by comparing the reconstructed consensus genomes to the original target genomes, defining three distinct performance zones based on empirical ANI thresholds (Supplementary Fig. S1):
\begin{itemize}
    \item \textbf{Unresolved Zone (>88.7\%):} At relatively high genetic similarity between target genomes in the mixture, the pipeline is unable to distinguish between the co-infecting strains, collapsing the mixture into a single consensus sequence per reference bin.
    \item \textbf{Ambiguous Zone (85.8\% – 88.7\%):} In this intermediate transition window, the pipeline successfully identifies multiple strains; however, the relatively high sequence similarity still permits reads to cross-map between references. This causes the reconstructed consensus sequences to accumulate mismatches in shared regions, resulting in lower fidelity to the original genomes despite initial separation by the assemblers.
    \item \textbf{High-Fidelity Resolution (< 85.8\%):} Below this similarity threshold, the strains are sufficiently divergent to minimize cross-mapping interference, allowing for the accurate reconstruction of distinct, high-quality consensus genomes.
\end{itemize}

\subsection*{S3. Influence of scaffolding reference}
The choice of scaffolding reference can considerably influence the completeness of the reconstructed genome. To quantify this effect, we analysed a subset of 7 Lassa virus samples from our own sequencing efforts (Lassa virus - Illumina Miseq data with UMIs). Samples were randomly selected based on their measured viral load during diagnostic testing, ensuring that the dataset encompassed different viral load categories (Supplementary Fig. S2). We compared consensus genomes generated using a gradient of scaffolding reference sequences ranging from 64\% to 100\% Average Nucleotide Identity (ANI) relative to the sample's baseline consensus.

To establish this gradient, we first generated a baseline consensus for each sample using the unclustered RVDB (U-RVDB) as the reference pool. We then retrieved all \textit{Mammarenavirus lassaense} (taxid: 3052310) sequences from NCBI and calculated their ANI against these baseline consensus genomes using MMseqs2 \citep{Steinegger2017-ci}. From this pool, we selected one scaffolding reference genome at every 2\% ANI interval for each sample (Supplementary Fig. S3). The Lassa virus samples exhibited differential coverage between segments (S segment mean ~231X, L segment mean ~68X; Supplementary Fig. S4), allowing us to evaluate scaffolding performance on both complete and fragmented assemblies.

By default, nf-core/viralmetagenome selects the cluster centroid as a scaffolding reference. To assess the reference influence, the pipeline’s code needed to be slightly modified to force the specified genomes as scaffolding reference. This version of the code repository is available in the branch \url{https://github.com/Joon-Klaps/viralmetagenome/tree/force-external-ref-centroid}.

The pipeline was then run with the following non-default parameters: \texttt{--with\_umi}, \texttt{--umi\_deduplicate both}, \texttt{--skip\_read\_classification}, \texttt{--skip\_precluster}, \texttt{--intermediate\_consensus\_caller ivar}, \texttt{--consensus\_caller ivar}.

To quantify the impact of the scaffolding reference on the final output, we compared the consensus genomes generated with the specifically supplied scaffold reference genomes against the baseline consensus genomes generated when the U-RVDB was supplied (without scaffold reference forcing, allowing automatic centroid and reference selection). We evaluated the sequence similarity as plotted in Figure 2, where we observe an almost linear correlation for the L segment when references were forced. For the S segment, this pattern is less apparent as \textit{de novo} assemblers could reconstruct near-complete genomes for most samples. In the L segment, where coverage was lower, nf-core/viralmetagenome excelled by connecting individual contigs to improve completeness.

% Start a new page
\newpage
\section*{Supplementary Figures}
\begin{figure}[htbp]
    \centering
    \includegraphics[width=\textwidth]{supplfig1}
    \caption{\textbf{Mixture resolution dynamics across a genetic diversity gradient.} This scatter plot characterises nf-core/viralmetagenome’s ability to resolve co-infections of varying genetic similarity. The Y-axis represents the Average Nucleotide Identity (ANI) between the two target genomes present in the simulated mixture (Target vs. Representative), while the X-axis indicates the ANI of the reconstructed consensus genomes compared to the constant target genome.}
    \label{fig:supp-fig1}
\end{figure}

\begin{figure}[htbp]
    \centering
    \includegraphics[width=0.5\textwidth]{supplfig2}
    \caption{\textbf{Distribution of diagnostic viral loads.} The distribution of diagnostic PCR Ct values for the available Lassa virus samples. Samples are classified into empirical categories (Super low, Low, Mid, High, Super high) based on their Ct values. Up to 3 random samples were selected from each category to ensure a diverse representation of viral loads.}
    \label{fig:supp-fig2}
\end{figure}

\begin{figure}[htbp]
    \centering
    \includegraphics[width=\textwidth]{supplfig3}
    \caption{\textbf{Sequencing coverage across Lassa virus baseline consensus genomes.} Consensus genomes used to determine the coverages are the baseline consensus genomes (generated with U-RVDB as a reference library supplied with \texttt{--reference\_pool}). A,C display per-base coverage depth (log\textsubscript{10} scale) along the genome. B,D display mean coverage per consensus. Overall, the S segment has a high read depth (median $\sim$231X), whereas the L segment shows lower coverage (median $\sim$68X) along with multiple areas of zero coverage, presenting a challenge for de novo assembly.}
    \label{fig:supp-fig3}
\end{figure}

\begin{figure}[htbp]
    \centering
    \includegraphics[width=\textwidth]{supplfig4}
    \caption{\textbf{Selection of scaffolding reference genomes across a similarity gradient.} To assess how reference divergence affects scaffolding, we compared all available NCBI \textit{Mammarenavirus lassaense} genomes against the baseline consensus (generated with U-RVDB as a reference library supplied with \texttt{--reference\_pool}) for 7 Lassa virus samples. The y-axis displays the similarity (\% ANI) of these candidate references towards the baseline consensus. For the benchmarking experiment, specific references were selected at $\sim$2\% ANI intervals, highlighted in full dots, to establish a gradient of similarity ranging from 64\% to 100\%.}
    \label{fig:supp-fig4}
\end{figure}

\newpage
\section*{Supplementary Tables}

\begin{longtable}{p{2cm}p{1.8cm}p{1.8cm}p{8.3cm}p{1cm}}
\caption{Overview of computational tools and methods used in the nf-core/viralmetagenome pipeline. The pipeline integrates multiple bioinformatics tools across different stages of viral metagenomic analysis, from preprocessing to consensus calling. Each tool was selected for its specific capabilities and performance characteristics relevant to viral sequence analysis.}
\label{table:S1} \\
\toprule
\textbf{Stage} & \textbf{Step} & \textbf{Tool} & \textbf{Explanation} & \textbf{Reference} \\
\midrule
\endfirsthead

\multicolumn{5}{c}{\textit{Table \ref{table:S1} continued from previous page}} \\
\toprule
\textbf{Stage} & \textbf{Step} & \textbf{Tool} & \textbf{Explanation} & \textbf{Reference} \\
\midrule
\endhead

\midrule
\multicolumn{5}{r}{\textit{Continued on next page}} \\
\endfoot

\bottomrule
\endlastfoot

Pre-processing & Read trimming & fastp & A high-performance preprocessing solution built in C++ that consolidates multiple quality control operations into a single workflow. Fastp delivers enhanced processing speed compared to traditional alternatives, like Trimmomatic, while handling adapter removal, quality assessment, base correction, read and filtering. Fastp also features automated adapter detection capabilities for various Illumina sequencing protocols. & \citep{Chen2018-tu} \\
\cmidrule{3-5}
 &  & Trimmomatic & A comprehensive preprocessing solution tailored for Illumina sequencing platforms with specialised paired-end read handling capabilities. Provides multiple pre-processing options, including adapter removal, quality-based trimming using sliding window approaches, and length-based filtering. & \citep{Bolger2014-si} \\
\cmidrule{2-5}
 & Complexity filtering & bbduk & A high-performance, multi-threaded tool designed to combine various data-quality-related operations, including trimming, filtering, and masking, into a single pass. It is particularly effective for removing host and other contaminant DNA using k-mer matching strategies, which is crucial for accurate metagenomic analysis. & \citep{Bushnell2014-qy} \\
\cmidrule{3-5}
 &  & PRINSEQ++ & A C++ implementation that significantly improves upon its predecessor, prinseq-lite.pl, in terms of computational efficiency. It offers extensive quality control features, including filtering by length, GC content, quality scores, N content, and sequence complexity (entropy/DUST score). & \citep{Cantu2019-vs} \\
\midrule

Metagenomic diversity & Taxonomy classification & Kraken2 & An advanced taxonomic classification system with an improved k-mer-based approach through optimised memory usage and enhanced processing speed. Incorporates protein-level search capabilities for improved detection of divergent viral sequences, with a library of pre-built RefSeq indexes on \url{https://benlangmead.github.io/aws-indexes/k2} & \citep{Wood2019-jl} \\
\cmidrule{3-5}
 & & Kaiju & A protein-level taxonomic classifier that employs exact matching algorithms on translated sequences using the Burrows-Wheeler transform. It can optionally allow amino acid substitutions and tends to have higher sensitivity and precision compared to k-mer-based classifiers (Kraken2). Kaiju is particularly effective for organisms with limited reference representation. A library of pre-built indexes, including RVDB, can be found on \url{https://bioinformatics-centre.github.io/kaiju/downloads.html}. & \citep{Menzel2016-tz} \\
\midrule

Assembly \& polishing & Assembly & SPAdes & A de Bruijn graph-based assembler designed to overcome challenges in complex sequencing data, such as non-uniform coverage (using multisized de Bruijn graphs) and variable insert sizes. The Spades suite contains multiple specialised modes rnaSPAdes, coronaSPAdes, metaSPAdes for extra flexibility. & \citep{Meleshko2021-gb} \\
\cmidrule{3-5}
 & & MEGAHIT & A succinct de Bruijn graph-based assembler specifically designed for large and complex metagenomic datasets. MEGAHIT has demonstrated the ability to generate larger assemblies with longer contig N50 and average contig length. Its relatively low memory requirements make it well-suited for handling large datasets. & \citep{Li2016-sd} \\
\cmidrule{3-5}
 & & Trinity & A powerful modular transcriptome reconstruction tool, capable of fully reconstructing a large fraction of transcripts, including alternative splice isoforms and transcripts from recently duplicated genes, even in the absence of a reference genome. It partitions sequence data into individual de Bruijn graphs, processing each independently to extract full-length isoforms and resolve paralogous genes. Trinity's sensitivity is comparable to methods relying on genome alignments. & \citep{Grabherr2011-ef} \\
\cmidrule{2-5}
 & Read alignment & BWAmem2 & An optimized version of the BWA-MEM algorithm enhanced with Intel-specific acceleration feature, providing approximately a twofold speedup in alignment throughput over BWA-MEM while ensuring identical SAM output. It is highly efficient for aligning short reads to large reference genomes. & \citep{Vasimuddin2019-rb} \\
\cmidrule{3-5}
 & & Bowtie2 & A fast and memory-efficient alignment tool for aligning sequencing reads to long reference sequences, such as mammalian genomes. Bowtie2 supports gapped, local, and paired-end alignment modes, combining high speed, sensitivity, and accuracy & \citep{Langmead2019-wx} \\
\midrule

Clustering & & CD-HIT-EST & A highly efficient and widely used program for clustering large sets of protein or nucleotide sequences. Implementation constraints limit processing to genomes under 10Mbp with minimum 0.8. The program employs a greedy incremental algorithm. & \citep{Li2006-nj} \\
\cmidrule{3-5}
 & & MMSeqs & A powerful software suite for fast and sensitive deep clustering and searching of large protein sequence sets. It is significantly faster than other tools like BLASTclust and can cluster large databases down to low sequence identity thresholds. MMseqs offers a variety of clustering modes for a lot of flexibility. & \citep{Steinegger2017-ci} \\
\cmidrule{3-5}
 & & VSEARCH & An open-source, versatile alternative for USEARCH. VSEARCH supports processing very large datasets, limited primarily by available memory. & \citep{Rognes2016-ju} \\
\cmidrule{3-5}
 & & Mash & A non-alignment based clustering technique that applies MinHash algorithms for rapid genome and metagenome distance estimation. It compresses large sequences into small, representative sketches, enabling ultra-fast estimations of global mutation distances. Mash creates compact sequence representations enabling ultra-fast similarity assessments and large-scale clustering operations, though accuracy decreases below 95\% genome similarity. & \citep{Ondov2019-bo} \\
\cmidrule{3-5}
 & & vRhyme & A machine learning-based binning solution specifically developed for viral genome recovery from metagenomic datasets. VRhyme addresses the unique challenges of viral sequences, such as the lack of universal marker genes. Addresses unique viral sequence characteristics through supervised learning approaches combined with coverage-based analysis across multiple samples. Generates high-quality viral bins with minimal computational overhead, though it may produce empty results when insufficient viral evidence is detected. & \citep{Kieft2022-km} \\
\midrule

Variant \& consensus calling & Variant \& consensus calling & iVar & A computational package specifically designed for viral amplicon-based sequencing, integrating functions like consensus and variant calling (including iSNVs and insertions/deletions). It is a key component of best-practice pipelines for reconstructing consensus genomes from viral sequencing data. & \citep{Grubaugh2019-xd} \\
\cmidrule{3-5}
 & & BCFtools & A robust suite of utilities for manipulating variant calls in VCF and BCF formats. Provides extensive functionality for variant calling using multiallelic models, data filtering, file merging, and consensus sequence generation through variant application. & \citep{Danecek2021-je} \\
\bottomrule
\end{longtable}
